\chapter{Graph Transformation and Compiler IR Rewriting}
\label{ch:compiler-graph-transformation}

A compiler graph contains more than expression children.
Value uses, control edges, block ownership, dominance, memory dependencies, and exception successors all constrain a valid transformation.
Deleting a node because its local pattern matches can leave a user pointing at nothing or a value definition on the wrong side of a branch.
Graph transformation therefore needs an interface between the matched region and the context that remains.

\section{Specify what the rule preserves}
\label{sec:ch87-interface}
A common algebraic description uses a span $L\leftarrow K\to R$.
The left graph L is matched, the interface K identifies preserved structure, and R describes the replacement including that interface.
The host graph supplies the surrounding context.
This description separates deleting structure, preserving structure, and adding structure instead of treating a rule as an unstructured before/after drawing.

\term{Double-pushout graph rewriting} formalizes deletion and gluing through categorical constructions under stated graph-category assumptions~\cite{graphtransformation-2006-fundamentals}.
For the ordinary injective setting, a central operational issue is the dangling condition: deleting a vertex must not leave an incident context edge without an endpoint.
Other frameworks, such as single-pushout approaches, can delete incident edges automatically, which gives different semantics.
A compiler cannot choose between those policies accidentally.

Figure~\ref{fig:ch87-interface} illustrates a producer-consumer replacement whose external input and output remain available.
The preserved boundary is a contract for clients outside the matched region.
For SSA IR, it includes type-compatible results and any control/memory token results, not only the main arithmetic value.
This chapter explains general graph-rewrite design, while the lab executes only a restricted deletion check with a supplied match.

\begin{figure}[htbp]
\centering
\begin{tikzpicture}[node distance=12mm]
\node[basic block] (i) {preserved external inputs};
\node[basic block,below=of i] (l) {matched internal region L};
\node[basic block,below=of l] (o) {preserved result interface K};
\draw[cfg edge] (i)--(l);\draw[cfg edge] (l)--(o);
\end{tikzpicture}
\caption{An IR rewrite must preserve the interface observed by the surrounding program.}
\label{fig:ch87-interface}
\end{figure}

\section{Check dangling references before deletion}
\label{sec:ch87-dangling}
Suppose vertices are $\{0,1\}$ with edge $0\to1$.
Deleting vertex one while preserving that edge is invalid.
If the rule explicitly deletes the edge too, a fuller deletion operation can be valid; if it redirects the edge to a replacement vertex, the use can remain satisfied.
The restricted companion deletes only isolated selected vertices, so it rejects any incident edge rather than silently choosing a redirection policy.

Table~\ref{tab:ch87-dangling} gives the executed cases.
The provided match is a set of vertex IDs, not a solved general subgraph-isomorphism problem.
General graph matching can be combinatorial; typing, fixed labels, bounded patterns, and IR-specific indexes often make practical compiler matching manageable.
Those engineering restrictions should be included before quoting a matching complexity bound.

\begin{table}[htbp]
\caption{Deleting a node and deleting its references are separate graph operations.}
\label{tab:ch87-dangling}
\centering
\begin{tabular}{lll}
\toprule
Edges & Deleted vertices & Restricted lab outcome \\
\midrule
$\{0\to1\}$ & $\{1\}$ & reject dangling edge \\
empty & $\{1\}$ & retain vertex zero \\
empty & unknown vertex two & reject unknown vertex \\
\bottomrule
\end{tabular}
\end{table}

Algorithm~\ref{alg:ch87-delete} checks that restricted contract transactionally.
With explicit sets, membership is expected constant time and checking every edge costs $O(|E|+|D|)$ for deleted set D, excluding output-set copying.
Adjacency indexes can check only incident edges but must themselves be maintained correctly.
The algorithm never mutates the original graph on failure.

\begin{algorithm}[htbp]
\caption{Delete supplied isolated vertices without inventing an edge-removal policy.}
\label{alg:ch87-delete}
\begin{pseudocode}
\Require vertex set $V$, valid edge set $E$, selected deletion set $D$
\Ensure graph with $D$ removed, or unchanged input and a diagnostic
\State require $D\subseteq V$
\ForAll{edges $(u,v)\in E$}
\If{$u\in D$ or $v\in D$} \State reject dangling incident edge \EndIf
\EndFor
\State return $(V\setminus D,E)$ as a new graph
\end{pseudocode}
\end{algorithm}

\section{Preserve SSA and control validity}
\label{sec:ch87-ssa}
Replacing every use of value v with value w requires compatible types and dominance of each new use.
An apparently equivalent value computed only in one branch cannot replace a value used after a join unless the necessary path conditions or merge structure are preserved.
For phi-like nodes, the use belongs to an incoming edge and its predecessor context, not simply the textual location of the merge instruction.
A CFG edit can also change the analyses used to justify the rewrite.

Figure~\ref{fig:ch87-dominance} shows a value produced on only one arm.
Moving its use to the join without a valid merge or dominating definition makes the IR ill-formed even before asking whether its numeric computation is equivalent.
Graph well-formedness and semantic preservation are different verification layers.
A graph-category gluing condition by itself does not establish dominance, memory ordering, or exception behavior.

\begin{figure}[htbp]
\centering
\begin{tikzpicture}[node distance=13mm and 17mm]
\node[basic block] (e) {branch};\node[basic block,below left=of e] (l) {define w};
\node[basic block,below right=of e] (r) {no w};\node[basic block,below=of e,yshift=-22mm] (j) {join uses v};
\draw[cfg edge] (e)--(l);\draw[cfg edge] (e)--(r);\draw[cfg edge] (l)--(j);\draw[cfg edge] (r)--(j);
\end{tikzpicture}
\caption{A branch-local definition cannot replace a join-visible value without a dominance argument.}
\label{fig:ch87-dominance}
\end{figure}

Memory operations add ordering constraints beyond value edges.
Replacing two loads by one requires an intervening-effect argument and compatible access semantics.
Deleting a call because its result is unused may remove I/O, exceptions, synchronization, or divergence.
The effect-aware IR chapter provides the summaries that such rules need; a graph matcher only locates the candidate region.

\section{Build and commit a replacement atomically}
\label{sec:ch87-transaction}
A replacement can fail halfway through allocation or invariant checking.
Publishing new uses before their definitions or deleting old nodes before all uses are redirected exposes an invalid intermediate graph to verifiers and concurrent readers.
Algorithm~\ref{alg:ch87-commit} uses a private replacement and explicit commit boundary.
The actual engine may use an undo log instead, but it must provide equivalent rollback and publication guarantees.

\begin{algorithm}[htbp]
\caption{Commit an IR graph rewrite only after structural and semantic checks succeed.}
\label{alg:ch87-commit}
\begin{pseudocode}
\Require valid IR, matched region and version, established rule premises
\Ensure valid committed replacement or unchanged visible IR
\State allocate replacement privately and connect its preserved interface
\State validate types, arities, dominance, control successors, and memory results
\State revalidate match and analysis versions at the commit boundary
\State redirect all selected uses according to the rule's sharing policy
\State publish replacement and retire unreachable old nodes safely
\State invalidate or update affected analyses and reverse-use indexes
\end{pseudocode}
\end{algorithm}

Local structural checks cost the size of the replacement and touched use lists.
Dominance or alias recomputation can involve the whole function, so a universal constant-time rewrite claim is usually false.
Incremental maintenance can reduce common-case cost but introduces its own consistency invariants.
An analysis invalidation table should say which edits preserve it, which update it, and which force recomputation.

Table~\ref{tab:ch87-invalidation} illustrates those dependencies.
A mere operator replacement preserving all control edges can still alter effects or use counts.
A CFG edit can invalidate dominance, loops, and phi-edge mappings even when no arithmetic opcode changes.
Recording only a generic modified flag makes both unnecessary recomputation and stale-analysis bugs harder to diagnose.

\begin{table}[htbp]
\caption{Different graph edits invalidate different compiler facts.}
\label{tab:ch87-invalidation}
\centering
\begin{tabular}{ll}
\toprule
Edit & Facts needing review \\
\midrule
replace value result & type compatibility, uses, value analyses \\
change memory operation & alias/effect summaries, memory dependencies \\
add or remove CFG edge & dominance, loops, merge inputs, reachability \\
split shared computation & use counts, cost, effect duplication \\
remove block & debug spans, exception edges, live metadata references \\
\bottomrule
\end{tabular}
\end{table}

\section{Parallelize independent matches only with conflicts defined}
\label{sec:ch87-parallel}
Two matches can be node-disjoint yet interfere through analysis premises or shared external uses.
For example, both can rely on a one-use condition that one rewrite changes, or both can redirect the same preserved interface.
A conflict relation should cover read/write sets of graph structure, analysis epochs, and ownership changes.
Optimistic parallel engines need validation and retry; serial engines still benefit from explicit conflict information for scheduling.

The lab does not implement parallel graph rewriting or a complete DPO framework.
Its supplied-match deletion fixture and shared-term update isolate two concrete hazards that larger implementations must handle.
The chapter's broader contracts connect those hazards to real compiler IR rather than implying that a few pointer edits constitute a general graph-transformation library.

\section{Use negative conditions and preserved identity deliberately}
\label{sec:ch87-negative}
Some graph rules require the absence of a structure, such as no additional uses of an intermediate result or no exceptional edge from a candidate region.
These are negative application conditions, not ordinary positive pattern edges.
A matcher finding all required nodes has not established an absence claim unless it also checks the complete relevant neighborhood.
An incomplete use list can therefore make a locally convincing graph rule unsound.

Preserved identity is equally important.
If a rule removes a wrapper but clients observe its address as an object identity, redirecting every use to its child changes an observable property even when payload values agree.
Pure compiler temporaries often have no observable identity, while heap allocations, synchronization objects, and foreign handles can.
The interface K should preserve the identities that the surrounding semantics can distinguish, not only the labels needed to redraw a similar graph.

Typed graph transformations can encode some constraints directly in node and edge sorts.
An edge representing a value use is different from an edge representing control succession or block ownership.
Requiring sort-preserving matches rules out many malformed replacements cheaply.
But type graphs generally do not express whole-path properties such as dominance or the absence of an intervening write.
Those remain semantic or analysis premises checked around the structural match.

\section{Plan CFG replacement as a multi-result operation}
\label{sec:ch87-control}
Consider replacing a conditional branch with a direct jump after proving its test constant.
Deleting the untaken edge changes the destination's predecessor set.
Any phi-like merge there must remove the corresponding incoming value, and newly unreachable blocks may have their own outgoing edges and metadata obligations.
The branch edit is thus not a one-edge operation when the IR invariants are included.
A transactional helper can gather the affected blocks, merge inputs, reachability changes, and analysis invalidations before publication.

Exception edges make the interface richer still.
Replacing a potentially trapping instruction with a total operation may legitimately remove an exception successor only after the no-trap premise is established under the language's error semantics.
Moving a trap into a different block can alter which side effects precede it.
An ordinary labeled-graph isomorphism check cannot establish that ordering equivalence.
The graph rule needs the same effect and control contract used by the semantic IR.

This leads to a useful implementation boundary: structural helpers maintain uses, ownership, predecessor/merge consistency, and version counters; semantic rule code supplies and checks the premises authorizing the change.
Neither layer should secretly assume the other has established everything.
The verifier can independently inspect the committed IR for structural correctness, while translation validation or rule proofs address the preserved program behavior.
Separating those responsibilities makes counterexamples much easier to localize.

\section{Exercises}
\label{sec:ch87-exercises}
\begin{enumerate}
\item\label{ex:ch87-dangle} Apply the restricted deletion algorithm to edge $0\to1$ and deletion set $\{1\}$, then to an empty edge set.
\item\label{ex:ch87-dominate} Explain why w in Figure~\ref{fig:ch87-dominance} cannot replace every use of v at the join.
\item\label{ex:ch87-rollback} A private replacement allocates successfully but fails a type check.
Which structures may be reclaimed and which visible structures must remain unchanged?
\item\label{ex:ch87-conflict} Two disjoint matches both rely on the same external node having one use.
Explain why graph-node disjointness alone does not prove parallel independence.
\end{enumerate}

\section{Further reading}
\label{sec:ch87-reading}
The term-rewriting literature supplies the local-rule foundation~\cite{ohlebusch-2002-rewriting}; the algebraic graph-transformation reference develops the broader gluing framework~\cite{graphtransformation-2006-fundamentals}.
The book's SSA, effect-IR, and validation chapters provide the additional compiler invariants that pure graph gluing does not express.
